\subsection{从题目要求到建模路线}

三个问题的难点不在于计算一个总分，而在于评价对象、任务需求和部署约束分别属于不同层次。若直接对公开成绩求平均，首先会混合不同测试尺度，其次会把同一 Benchmark Family 中的重复信息计算多次，最后还可能把“未公开”错误解释为“能力最低”。因此，问题一需要先处理数据语义，再建立可比的模型间关系。

问题一得到的综合排序仍不能直接用于所有任务。科研长文本分析需要在复杂推理、事实可靠性和长上下文之间形成协同；代码开发更依赖复杂推理与代码工程能力的共同支撑；日常通用对话则允许更广泛的能力组合。若采用线性加权，某一项能力的不足可以被其他高分能力完全补偿，这与存在关键短板的任务特征不符。因此，问题二采用场景权重与 CES 聚合，将能力结构转译为任务效用。

场景效用解决了适配性问题，但仍没有回答“预算允许时应该选谁”。价格本身也不是任务成本：输入/输出 Token 数量和调用次数会改变各模型的实际支出，代码任务的输出负载较高时，输出单价差异尤其容易被放大。问题三因此保留场景效用，只增加工作负载成本，并以 Pareto 前沿和预算区间表达部署选择。

\subsection{全文建模框架}

图~\ref{fig:framework}将上述逻辑归纳为三层。第一层从公开评测得到模型的基础能力；第二层把能力映射为科研、日常和代码三类任务的效用；第三层在成本和预算约束下确定可行的优选节点。图中第一层的缺失与冗余处理保证“有什么数据”不被误读，第二层的场景权重回答“任务需要什么”，第三层的 Pareto 与预算区间回答“在现实约束下选谁”。因此，三问不是三个相互独立的排名，而是由能力识别逐步走向任务决策。

\begin{figure}[htbp]
\centering
\includegraphics[width=\textwidth]{fig_framework_user.jpg}
\caption{全文建模框架。第一层统一公开评测并识别模型能力，第二层将五维能力映射为场景需求，第三层在工作负载成本和预算约束下形成部署决策。}
\label{fig:framework}
\end{figure}
